% =============================================================================
% Abstract
% paper/abstract.tex
% =============================================================================

\begin{abstract}
Danielson, Satishchandran, and Wald showed that a body in a spatial quantum
superposition near a Killing horizon inevitably decoheres: soft radiation
crossing the horizon carries which-path information that is permanently lost
to any local observer.  The decoherence rate is set by the Hawking temperature,
tying the effect to the thermal character of the horizon.  We extend this
result to cosmological spacetimes that possess a future causal horizon but no
Killing symmetry and no thermal structure.  Working in the family of spatially
flat Friedmann-Lema\^{i}tre-Robertson-Walker (FLRW) spacetimes with power-law
scale factor $a(\tau) \propto \tau^{p}$, $p > 1$, we exploit the conformal
invariance of the massless scalar and Maxwell fields to reduce the decoherence
integral to a closed-form expression.  The decoherence number for a conformally
coupled scalar field is
\[
  \mathcal{N}(T)
  =
  \frac{q^{2}d^{2}p^{3}}{48\pi^{3}}
  \cdot
  \frac{2T_{0}^{2} + \widetilde{T}(2 + \widetilde{T})}
       {T_{0}^{2}(1 + \widetilde{T})^{2}}
  \cdot
  \ln(1 + \widetilde{T}) \,,
\]
where $q$ is the charge, $d$ the superposition separation, $T$ the
superposition duration, $T_{0}$ a proper-time scale set by the initial
conditions, and $\widetilde{T} = T/T_{0}$.  For large $T$,
$\mathcal{N} \sim \ln T$: decoherence grows logarithmically without bound,
confirming that the causal horizon produces \emph{inevitable} decoherence
even without any thermal bath.  Specializing to de~Sitter spacetime recovers
exactly the known result $\mathcal{N} = q^{2}d^{2}H^{3}T/(96\pi^{2})$,
with its linear growth and Gibbons-Hawking temperature dependence.  When
$p = 1$ (no future horizon), $\mathcal{N}$ saturates to a finite constant,
sharply demonstrating that the horizon is the essential mechanism.  For the
electromagnetic field, conformal invariance gives
$\mathcal{N}_{\rm EM} = 2\,\mathcal{N}_{\rm scalar}$ universally, counting
the two transverse photon polarizations.
\end{abstract}
